\section{Dimensional and scaling inference: proofs}
\label{sec:annex-p15-dimensional-and-scaling-inference}

\begin{annexitem}{annex:p15.dimensional.rank-nullity}{thm:p15-dimensional-rank-nullity}{Dimension vectors and the free exponents}
\begin{proof}
Under the stated unit change,
\[
 (x')^a=x^a\exp(a^{\mathsf T}D^{\mathsf T}s)
        =x^a\exp((Da)^{\mathsf T}s).
\]
This equals \(x^a\) for every \(s\) exactly when \(Da=0\):
if \(Da\ne0\), take \(s=Da\).

Elementary row operations replace a homogeneous system by an equivalent
one. In row-echelon form there are \(r\) pivot variables and \(n-r\)
free variables. Each pivot is uniquely determined by the free variables.
Setting the free variables successively to their coordinate vectors gives
\(n-r\) independent solutions, and every solution is their linear
combination. This proves the dimension and supplies a constructive basis.

If \(a=\sum_jc_jb_j\), positivity and the exponential definition of powers
give \(x^a=\prod_j(x^{b_j})^{c_j}\). Replacing dimension coordinates
by an invertible matrix \(A\) replaces \(D\) with \(AD\);
\(ADa=0\) is equivalent to \(Da=0\).
The displayed weight example has independent first and second rows and
third row \(-2\) times the first, so its rank is two and its nullspace
is exactly the stated one-dimensional set. The acceleration matrix has
two independent rows; direct substitution gives its stated basis.
\end{proof}
\end{annexitem}

\begin{annexitem}{annex:p15.dimensional.pi-factorization}{thm:p15-dimensional-pi-factorization}{Why the products describe complete unit-change orbits}
\begin{proof}
Put \(z=\log x\). Unit changes act by \(z\mapsto z+D^{\mathsf T}s\).
Because the rows of \(B\) are in \(\ker D\), \(BD^{\mathsf T}=0\).
The row and column ranks of a finite matrix agree, as follows from
row reduction and the corresponding pivot columns. Thus
\(\operatorname{im}D^{\mathsf T}\) has dimension \(r\).
The matrix \(B\) has rank \(n-r\), so
\(\ker B\) has dimension \(r\), by the preceding theorem.
The inclusion \(\operatorname{im}D^{\mathsf T}\subseteq\ker B\)
is therefore equality.

Consequently \(Bz=Bz'\) if and only if \(z'-z=D^{\mathsf T}s\)
for some \(s\). Exponentiation is injective on each real coordinate,
so the fibers of \(\pi\) are exactly the unit-change orbits.
An invariant \(F\) is constant on each fiber and hence defines a
well-defined \(\Psi\) on the product values. Conversely any such
factorization is constant on these fibers and is invariant.
Since \(B\) has full row rank, every real product-log vector is attained;
thus \(\pi\) maps onto \((0,\infty)^{n-r}\).
When \(n-r=0\), this target is the one-point empty-product space.
Applying the same argument to the indicator of an invariant set proves
the final assertion.
\end{proof}
\end{annexitem}

\begin{annexitem}{annex:p15.dimensional.covariant-output}{thm:p15-dimensional-covariant-output}{The remaining function in a dimensional output formula}
\begin{proof}
If \(Da=w\), the positive monomial \(x^a\) has precisely the output's
multiplicative transformation factor. The ratio \(f(x)/x^a\) is therefore
invariant. Apply \cref{thm:p15-dimensional-pi-factorization}.

If \(w\) is outside the column space of \(D\), finite-dimensional
orthogonal decomposition gives a vector \(s\) orthogonal to that space
with \(w^{\mathsf T}s\ne0\). One can obtain it by subtracting from \(w\)
its projection onto an orthonormal basis of the column space.
Then \(D^{\mathsf T}s=0\), so this unit change fixes every input value,
whereas covariance requires \(f(x)=e^{w^{\mathsf T}s}f(x)\).
The multiplier differs from one, forcing \(f(x)=0\) for every \(x\).
The required projection uses only a finite independent basis and
Gram--Schmidt orthogonalization.
\end{proof}
\end{annexitem}

\begin{annexitem}{annex:p15.dimensional.scaling-calculus}{thm:p15-dimensional-scaling-calculus}{Substitution and addition of exponent contributions}
\begin{proof}
For positive \(x,A\),
\((Ax^\alpha)^\beta=A^\beta x^{\alpha\beta}\), by the logarithmic
definition of real powers. This proves the chain rule for the first pair
of monomials. In the general finite product, substitute each expression
for \(u_j\) and collect the coefficients of each \(\log x_i\):
\[
 \log y=\log K+\sum_jp_j\log A_j+
                         \sum_i\left(\sum_jp_j a_{ji}\right)\log x_i.
\]
Exponentiating gives the claimed formula.
For contrast, \(Ax^\alpha+Bx^\beta\) is a sum, whose logarithm is
not the sum of the two logarithms. When \(\alpha\ne\beta\), it need
not be a single power law.
\end{proof}
\end{annexitem}

\begin{annexitem}{annex:p15.dimensional.log-slope}{thm:p15-dimensional-log-slope}{Constant elasticity and a doubling counterexample}
\begin{proof}
For \(y=Kx^\alpha\),
\(\eta=\log(Kx_0^\alpha/y_0)+\alpha\xi\), giving a straight line.
For positive differentiable \(f\), the chain rule gives
\(d\log f(x)/d\log x=xf'(x)/f(x)\).
If this equals \(\alpha\) on an interval, the derivative of
\(\log f(x)-\alpha\log x\) is zero. The mean-value theorem makes
that expression constant, and exponentiation yields a power law.

For the body's counterexample, replacing \(x\) with \(2x\) increases
\(\log_2x\) by one, leaving the sine factor unchanged.
But the logarithmic derivative is
\[
 \alpha+\frac{2\pi\varepsilon}{\log2}
                \cos(2\pi\log_2x),
\]
which is not constant when \(\varepsilon\ne0\).
Thus a doubling experiment alone is not the continuous scaling premise.
\end{proof}
\end{annexitem}

\begin{annexitem}{annex:p15.dimensional.mechanical-models}{thm:p15-dimensional-mechanical-models}{Geometric scaling and the stipulated equations of motion}
\begin{proof}
Let \(\lambda_2\) denote planar Lebesgue measure.
For a planar rectangle, multiplying both lengths by \(s\) multiplies
area by \(s^2\). The measures \(E\mapsto s^{-2}\lambda_2(sE)\) and
\(E\mapsto\lambda_2(E)\) therefore agree on rectangles and are
sigma-finite. Uniqueness of the product measure extends their agreement
to Borel sets. Taking an area ratio of \(1/2\) gives length ratio
\(s=1/\sqrt2\).

Uniform circular motion is parametrized by
\(r(\cos\omega t,\sin\omega t)\). Its velocity has constant magnitude
\(r\omega\), its acceleration is \(-\omega^2\) times its position,
and a full revolution takes \(2\pi/\omega\).
Equating the inward acceleration \(v^2/r\) with the stipulated
\(GM/r^2\) gives \(v^2=GM/r\) and period
\(2\pi r/v=2\pi\sqrt{r^3/(GM)}\).

For the projectile, with origin at launch,
\[
 x(t)=v\cos\theta\,t,\qquad
 y(t)=v\sin\theta\,t-\tfrac12gt^2.
\]
These functions satisfy the specified constant acceleration and initial
velocity. The nonzero return time to \(y=0\) is \(2v\sin\theta/g\).
Substitution in \(x(t)\) gives
\(2v^2\sin\theta\cos\theta/g=v^2\sin(2\theta)/g\).

Put \(\omega=\sqrt{k/m}\) for the linear oscillator.
Every initial position and velocity are matched by
\(A\cos\omega t+B\sin\omega t\).
For the difference \(z\) between this function and any other solution
with the same initial values, differentiation gives
\(\frac{d}{dt}(\dot z^2+\omega^2z^2)=0\).
Its initial value is zero, so \(z=0\). A nonzero solution is a sinusoid
with amplitude \(\sqrt{A^2+B^2}>0\); its least positive period is
\(2\pi/\omega\). This proves the period within the linear equation,
not exact amplitude independence for a nonlinear pendulum.
\end{proof}
\end{annexitem}

\begin{annexitem}{annex:p15.dimensional.competing-powers}{thm:p15-dimensional-competing-powers}{Balancing weighted contributions at the unique optimum}
\begin{proof}
Differentiation gives
\(E'(v)=-pAv^{-p-1}+qBv^{q-1}\).
Multiplication by the positive \(v^{p+1}\) leaves its sign equal to
that of \(-pA+qBv^{p+q}\), a strictly increasing expression with a
unique zero at the stated \(v_*\).
The derivative is negative before and positive after this point.
Also \(E(v)\to\infty\) as \(v\downarrow0\) or \(v\to\infty\).
Thus this is the unique global minimum.
The derivative equation gives the balance of weighted contributions.

For \(p=q=2\), the two contributions equal \(\sqrt{AB}\), giving
\(E_{\min}=2\sqrt{AB}\) and \(v_E=(A/B)^{1/4}\).
For \(P(v)=A/v+Bv^3\), the same argument with exponents \(1,3\)
gives \(v_P=(A/(3B))^{1/4}=3^{-1/4}v_E\).

Substituting the body's \(A,B\) gives energy per distance
\(2\sqrt C\,mg\) and
\(v_E=(mg)^{1/2}C^{-1/4}\rho^{-1/2}L^{-1}\).
Under fixed shape and mass density \(m\propto L^3\), this becomes
\(v_E\propto m^{1/6}\rho^{-1/2}\), with \(C,g\) fixed.
Changing the objective from energy per distance to power changes the
optimum; equality of the two energy contributions was specific to
the equal exponents.
\end{proof}
\end{annexitem}

\begin{annexitem}{annex:p15.dimensional.birthday}{thm:p15-dimensional-birthday}{Exact occupancy probabilities and the logarithmic approximation}
\begin{proof}
Conditional on the first \(j\) birthdays being distinct, exactly \(M-j\)
of the \(M\) categories avoid a match for the next independent birthday.
Multiplying these conditional probabilities gives the product.
For \(n>M\), the pigeonhole principle forces a match.

For \(0\leq t\leq a<1\),
\[
 -\log(1-t)-t=\int_0^t\frac{s}{1-s}\,ds,
 \qquad
 0\leq-\log(1-t)-t\leq\frac{t^2}{2(1-a)}.
\]
Apply this to \(t=j/M\) and sum. The first-order sum is
\(S_1=n(n-1)/(2M)\); the squared sum is
\(n(n-1)(2n-1)/(6M^2)\). This proves both log bounds.
The remainder is small, for example, when \(n=O(\sqrt M)\) as
\(M\to\infty\).

For occupancies \(n_1,\ldots,n_M\) with sum \(n\), the number of
assignments of labelled persons to those counts is
\(n!/\prod_b n_b!\). Sum this over \(0\leq n_b<k\) and divide by
\(M^n\). The resulting finite sum is exactly the displayed coefficient.
It can be evaluated with integers by
\[
 A_0(0)=1,\quad A_0(t)=0\ (t>0),\qquad
 A_b(t)=\sum_{j=0}^{\min(k-1,t)}\binom tj A_{b-1}(t-j).
\]
Here \(A_b(t)\) counts assignments to \(b\) labelled categories without a
\(k\)-fold occupancy. Thus the noncollision probability is \(A_M(n)/M^n\).
This finite recurrence gives the numerical pair and triple comparisons in
the body. Collision probability is nondecreasing as persons are added,
so the comparisons on adjacent integers establish the first threshold.

Finally let \(N\) count \(k\)-subsets whose birthdays coincide.
Each such subset has coincidence probability \(M^{-(k-1)}\);
linearity gives \(EN=\binom nk/M^{k-1}\).
The inequality \(P(N\geq1)\leq EN\) is the union bound.
No independence of the subset events has been assumed or inferred.
\end{proof}
\end{annexitem}

\begin{annexitem}{annex:p15.dimensional.product-bounds}{thm:p15-dimensional-product-bounds}{Endpoint and logarithmic bounds for an estimate tree}
\begin{proof}
For positive \(x_i\), \(x_i^{a_i}\) is increasing, constant or decreasing
according as \(a_i\) is positive, zero or negative.
Choose each endpoint accordingly and multiply. Every chosen endpoint
combination is in the rectangular domain, proving the minima and maxima
are attained there.
If \(x_i=\widehat x_i e^{\delta_i}\), then
\[
 \log\frac{K\prod_i x_i^{a_i}}{K\prod_i\widehat x_i^{a_i}}
             =\sum_i a_i\delta_i.
\]
The triangle inequality gives the claimed bound.
No sampling model or distribution was used. Restricting feasible input
combinations to a subset keeps these as bounds but may remove attainment.

For the source's tree convention, each internal product node substitutes
monomials from its descendants. Repeated use of
\cref{thm:p15-dimensional-scaling-calculus} gives the product of leaf
values with the accumulated path exponents. A tree containing a sum must
retain that sum; it is not changed into a monomial by the notation.
\end{proof}
\end{annexitem}

\begin{annexitem}{annex:p15.dimensional.source-extensions}{thm:p15-dimensional-source-extensions}{Recurrences and the remaining conditional scaling calculations}
\begin{proof}
For \(n=2^h\), induction in \(h\) gives
\(C(n)=n\log_2n-n+1\): substitute the asserted expression for \(n/2\)
into the recurrence. Likewise \(M(2^h)=3^h=(2^h)^{\log_2 3}\);
replacing three with four gives \(n^2\).
The three-product algebra behind Karatsuba is
\[
 (aB+b)(cB+d)
 =acB^2+\{(a+b)(c+d)-ac-bd\}B+bd.
\]
This identity explains the recursive-product model.
Addition, carries, uneven digit lengths and base-case choices must
be accounted for in a concrete implementation's running time.

For uniform density \(\rho\), the stipulated enclosed mass at \(r<R\)
is \(4\pi\rho r^3/3\). Substituting it into \(Gm(r)/r^2\) gives
\(g(r)=4\pi G\rho r/3\). Outside the sphere, the enclosed mass is
the constant \(4\pi\rho R^3/3\), giving inverse-square dependence.

For the quartic oscillator, multiplication of
\(m\ddot x=-4Cx^3\) by \(\dot x\) shows that
\(\frac12m\dot x^2+Cx^4\) is constant. At a turning point it is
\(Cx_0^4\). On a monotone quarter-cycle,
\[
 |dt/dx|=\sqrt{\frac m{2C}}\frac1{\sqrt{x_0^4-x^4}}.
\]
Substitute \(x=x_0u\) and multiply the travel time from zero to \(x_0\)
by four. The even potential and restoring equation give identical
quarter-cycle times. The integral is finite, since
\(1-u^4\geq1-u\) on \([0,1]\), so its integrand is bounded above by
\((1-u)^{-1/2}\). The nonzero restoring acceleration at a turning point
selects the returning branch; arbitrary waiting at a turning point is not
a solution of the stated second-order equation.

In the rolling model, torque balance gives \(f=Ia/r^2\).
Substituting this into translation balance yields
\(mg\sin\theta=a(m+I/r^2)\). Divide by the strictly positive
coefficient to obtain the claimed formula.
\end{proof}
\end{annexitem}
